\documentclass[12pt, a4paper]{article}

% ── Packages ──────────────────────────────────────────────────────────────────
\usepackage[margin=1in]{geometry}
\usepackage[T1]{fontenc}
\usepackage[utf8]{inputenc}
\usepackage{lmodern}
\usepackage{microtype}
\usepackage{xcolor}
\usepackage{titlesec}
\usepackage{booktabs}
\usepackage{tabularx}
\usepackage{array}
\usepackage{colortbl}
\usepackage{enumitem}
\usepackage{listings}
\usepackage{hyperref}
\usepackage{fancyhdr}
\usepackage{graphicx}
\usepackage{tcolorbox}
\usepackage{setspace}
\usepackage{mdframed}
\usepackage{amsmath}
\tcbuselibrary{skins, breakable}

% ── Colors ────────────────────────────────────────────────────────────────────
\definecolor{primaryblue}{HTML}{1A73E8}
\definecolor{darkblue}{HTML}{0D47A1}
\definecolor{lightgray}{HTML}{F5F5F5}
\definecolor{codebg}{HTML}{1E1E2E}
\definecolor{codetext}{HTML}{CDD6F4}
\definecolor{accent}{HTML}{E53935}
\definecolor{mysqlorange}{HTML}{E8710A}
\definecolor{tableheader}{HTML}{1A73E8}

% ── Hyperref Setup ────────────────────────────────────────────────────────────
\hypersetup{
    colorlinks=true,
    linkcolor=primaryblue,
    urlcolor=primaryblue,
    pdftitle={GitDB: Version Control System for MySQL Databases},
    pdfauthor={Group Project - DBMS}
}

% ── Section Formatting ────────────────────────────────────────────────────────
\titleformat{\section}
  {\large\bfseries\color{darkblue}}
  {\thesection.}{0.5em}{}[\color{primaryblue}\titlerule]

\titleformat{\subsection}
  {\normalsize\bfseries\color{primaryblue}}
  {\thesubsection.}{0.5em}{}

\titleformat{\subsubsection}
  {\normalsize\bfseries\color{primaryblue!80}}
  {\thesubsubsection.}{0.5em}{}

\titlespacing*{\section}{0pt}{1.2em}{0.6em}
\titlespacing*{\subsection}{0pt}{0.8em}{0.4em}
\titlespacing*{\subsubsection}{0pt}{0.6em}{0.3em}

% ── Code Listing Style ────────────────────────────────────────────────────────
\lstset{
    backgroundcolor=\color{codebg},
    basicstyle=\ttfamily\small\color{codetext},
    keywordstyle=\color{cyan}\bfseries,
    commentstyle=\color{gray}\itshape,
    stringstyle=\color{yellow},
    breaklines=true,
    frame=none,
    xleftmargin=1em,
    xrightmargin=1em,
    aboveskip=0.5em,
    belowskip=0.5em,
    showstringspaces=false,
    numbers=none,
}

% ── Custom Boxes ──────────────────────────────────────────────────────────────
\tcbset{
    commandbox/.style={
        colback=codebg,
        colframe=primaryblue,
        coltext=codetext,
        fonttitle=\bfseries\color{white},
        attach boxed title to top left={yshift=-2mm, xshift=4mm},
        boxed title style={colback=primaryblue, colframe=primaryblue},
        arc=4pt,
        boxrule=0.8pt,
        left=6pt, right=6pt, top=8pt, bottom=6pt,
        breakable
    },
    infobox/.style={
        colback=lightgray,
        colframe=primaryblue,
        arc=4pt,
        boxrule=0.8pt,
        left=8pt, right=8pt, top=6pt, bottom=6pt
    },
    warningbox/.style={
        colback=mysqlorange!8,
        colframe=mysqlorange,
        arc=4pt,
        boxrule=0.8pt,
        left=8pt, right=8pt, top=6pt, bottom=6pt
    }
}

% ── Header / Footer ───────────────────────────────────────────────────────────
\pagestyle{fancy}
\fancyhf{}
\fancyhead[L]{\small\textcolor{primaryblue}{\textbf{GitDB}} — Version Control for MySQL}
\fancyhead[R]{\small\textcolor{gray}{DBMS Mini Project}}
\fancyfoot[C]{\small\textcolor{gray}{\thepage}}
\renewcommand{\headrulewidth}{0.4pt}
\renewcommand{\headrule}{\hbox to\headwidth{\color{primaryblue}\leaders\hrule height \headrulewidth\hfill}}

% ── Document ──────────────────────────────────────────────────────────────────
\begin{document}

% ── Title Page ────────────────────────────────────────────────────────────────
\begin{titlepage}
    \centering
    \vspace*{1.5cm}

    {\Huge\bfseries\color{darkblue} GitDB}\\[0.4em]
    {\Large\color{primaryblue} Version Control System for MySQL Databases}\\[0.2em]
    {\small\color{mysqlorange}\textit{Git for MySQL}}\\[0.2em]
    {\small\color{gray}\rule{8cm}{0.5pt}}

    \vspace{1.5cm}

    \begin{tcolorbox}[infobox, width=0.85\textwidth]
        \centering
        {\normalsize\textit{A MySQL-native, state-based database versioning tool with Git-like semantics,\\
        providing commit, diff, checkout, and restore operations\\
        for MySQL databases --- with all metadata stored in a dedicated \texttt{gitdb\_meta} database.}}
    \end{tcolorbox}

    \vspace{1.5cm}

    {\large\textbf{Institution:} PICT, Pune}\\[0.5em]
    {\large\textbf{Subject:} DISL}\\[0.5em]
    {\large\textbf{Project Type:} Mini Project}\\[0.5em]
    {\large\textbf{Class:} SY-11 \quad|\quad \textbf{Lab Batch:} F-11}\\[0.5em]
    {\large\textbf{Group Size:} 4 Members}\\[1em]

    \vspace{0.5cm}

    \textbf{Team Members:}\\
    \vspace{0.5em}
    \renewcommand{\arraystretch}{1.3}
    \begin{tabular}{>{\centering\arraybackslash}p{1cm} >{\arraybackslash}p{5cm} >{\centering\arraybackslash}p{2.5cm}}
    \toprule
    \textbf{Sr.} & \textbf{Name} & \textbf{Roll No.} \\
    \midrule
    1 & Aditya Malode     & 23325 \\
    2 & Rohan Natekar     & 23330 \\
    3 & Aarush Oak        & 23332 \\
    4 & Prathamesh Kinagi & 23340 \\
    \bottomrule
    \end{tabular}

    \vspace{2.5cm}
    {\small\color{gray} Academic Year 2025--26}
    \vfill
    {\small\color{gray} Compiled on \today}
\end{titlepage}

% ── Abstract ──────────────────────────────────────────────────────────────────
\section*{Abstract}
\addcontentsline{toc}{section}{Abstract}

\begin{onehalfspacing}

Version control has become a foundational pillar of modern software engineering.
Systems like Git allow developers to snapshot code, inspect historical changes,
revert to any prior state, and collaborate across teams with full auditability.
Yet despite MySQL databases being a critical component of almost every production
software system, they have never enjoyed an equivalent first-class versioning
mechanism. Schema changes are typically managed through manually written migration
scripts that are fragile, difficult to audit, and entirely disconnected from the
data they govern. Data rollbacks, when needed, require either expensive full
database dumps or complex reconstruction --- neither of which resembles the
intuitive, commit-and-revert workflow developers take for granted with source code.

\textbf{GitDB} addresses this gap directly. It is a \textbf{MySQL-native},
state-based version control system designed to bring Git-like semantics ---
commit, log, diff, and checkout --- to MySQL database management without requiring
developers to manually author SQL migration scripts. The central idea is simple:
instead of tracking \textit{instructions} for how to reach a state, GitDB tracks
the \textit{state itself}.

At every commit, a complete snapshot of the target MySQL database's schema (DDL)
and data is captured and persisted into a dedicated metadata database named
\texttt{gitdb\_meta}, which runs as a separate MySQL database on the same server.
This clean separation mirrors the way Git's \texttt{.git/} directory is isolated
from the source files it tracks --- the versioning layer never pollutes the
database being versioned. The \texttt{gitdb\_meta} database is built around
\textbf{four entities}: \texttt{user} (authentication and authorship),
\texttt{repository} (target database connection config and HEAD pointer),
\texttt{commits} (hash, message, author, timestamp, parent hash), and
\texttt{snapshots} (serialised schema and data per table per commit). The HEAD
pointer is stored directly as \texttt{current\_hash} inside \texttt{repository},
keeping the design normalised and eliminating the need for a separate head table.

Each commit is identified by a SHA-256 hash computed from its content
serialisation and the hash of its parent commit, forming an immutable,
tamper-evident \textbf{Merkle chain}. Since GitDB does not support branching,
commit history is always a \textbf{linear linked list} --- each commit has at
most one parent and at most one child. The system's core technical contribution
is its \textbf{diff engine}, which takes two commit snapshots from
\texttt{gitdb\_meta} and automatically generates the MySQL SQL necessary to
transition the live database from one state to the other. Schema diffing is
performed at the Abstract Syntax Tree (AST) level using \texttt{sqlglot},
correctly identifying added tables, dropped tables, and altered columns. Data
diffing uses a primary-key-anchored algorithm operating in $O(n)$ time,
producing \texttt{INSERT}, \texttt{UPDATE}, and \texttt{DELETE} statements for
rows that are new, changed, or removed.

A key engineering challenge specific to MySQL is that DDL statements
(\texttt{ALTER TABLE}, \texttt{DROP TABLE}, \texttt{CREATE TABLE}) issue implicit
commits and cannot be rolled back in a standard transaction. GitDB handles this
through a \textbf{two-phase checkout strategy}: schema changes are applied and
validated first; data patches are then applied inside an explicit transaction
with rollback on failure. If schema application fails, the pre-checkout snapshot
in \texttt{gitdb\_meta} is used to restore the prior state, guaranteeing the
database is never left inconsistent.

The system is delivered through two complementary interfaces. A
\textbf{command-line interface} built with \texttt{click} and \texttt{rich}
provides Git-familiar commands: \texttt{gitdb init}, \texttt{gitdb commit},
\texttt{gitdb log}, \texttt{gitdb diff}, \texttt{gitdb checkout}, and
\texttt{gitdb status}. A \textbf{web-based user interface} built in React~18
with Tailwind CSS offers a visual commit graph rendered using \texttt{react-flow}
and a side-by-side diff viewer via \texttt{@git-diff-view/react}, communicating
with a lightweight Flask REST API.

Academically, GitDB demonstrates MySQL schema introspection via
\texttt{information\_schema}, the practical implications of MySQL's DDL
auto-commit behaviour on transaction design, structural DDL comparison at the
AST level, cryptographic integrity via Merkle hashing, and the foundational role
of primary keys as stable row identifiers --- all demonstrated through working
code against the world's most widely deployed open-source relational database.

\end{onehalfspacing}

\vspace{1.5em}

\begin{tcolorbox}[infobox]
\textbf{Keywords:} MySQL, Version Control, State-Based Snapshotting, Schema
Diffing, SQL Patch Generation, \texttt{gitdb\_meta}, Merkle Hash Chain, Linear
Commit History, DDL Auto-Commit, Two-Phase Checkout, Python, React, Flask
\end{tcolorbox}

% ── Section 1: Introduction ───────────────────────────────────────────────────
\newpage
\section{Introduction and Motivation}

MySQL is the world's most widely deployed open-source relational database,
powering applications from small web projects to large-scale production systems.
Yet despite this ubiquity, MySQL development workflows lack a native, intuitive
mechanism for versioning both schema and data together. Schema changes are managed
through ad-hoc \texttt{ALTER TABLE} scripts, Flyway migrations, or Liquibase
changelogs --- none of which version data alongside schema, and none of which
offer Git's commit-and-revert ergonomics.

GitDB fundamentally reimagines MySQL database version control by adopting a
\textbf{state-based approach} instead of an instruction-based one. Rather than
tracking how to reach a state, GitDB tracks the state itself through complete
snapshots stored in a dedicated \texttt{gitdb\_meta} MySQL database.

\vspace{0.5em}
\textbf{Problems with the status quo:}

\begin{itemize}[leftmargin=1.5em, itemsep=0.3em]
    \item \textbf{Fragility:} Migration scripts can fail mid-execution, leaving
    MySQL databases in undefined states
    \item \textbf{Audit Difficulty:} Understanding current schema requires
    mentally replaying all migrations in order
    \item \textbf{Data Disconnection:} Schema migrations are entirely separate
    from data versioning
    \item \textbf{No Easy Rollback:} Reverting requires writing additional
    down-migrations or restoring from \texttt{mysqldump} backups
    \item \textbf{MySQL DDL Complexity:} MySQL DDL statements auto-commit,
    making transactional schema management non-trivial
\end{itemize}

% ── Section 2: Core Features ──────────────────────────────────────────────────
\section{Core Features and Objectives}

\subsection{MySQL-Native State-Based Snapshotting}

At every commit, GitDB captures a complete snapshot of the target MySQL database
using \texttt{information\_schema} and \texttt{SHOW CREATE TABLE} for schema, and
\texttt{SELECT * ... ORDER BY \{pk\_columns\}} for data. Both are serialised as JSON
and written into the \texttt{snapshots} table of \texttt{gitdb\_meta}.

\subsection{The \texttt{gitdb\_meta} Database}

All versioning metadata lives in a dedicated MySQL database named
\texttt{gitdb\_meta} on the same MySQL server instance. This is analogous to
Git's \texttt{.git/} directory --- the versioning layer is fully isolated from
the target database being versioned.

The metadata schema consists of \textbf{four entities}. The HEAD pointer ---
which tracks the current commit for each repository --- is stored as the
\texttt{current\_hash} column directly inside the \texttt{repository} table,
eliminating the need for a separate \texttt{head} table and keeping the design
cleanly normalised.

\vspace{0.5em}
\renewcommand{\arraystretch}{1.4}
\begin{tabularx}{\textwidth}{>{\bfseries}p{3cm} >{\ttfamily}p{4.5cm} X}
\rowcolor{tableheader}
\textcolor{white}{\normalfont\bfseries Entity} &
\textcolor{white}{\normalfont\bfseries Key Columns} &
\textcolor{white}{\normalfont\bfseries Purpose} \\
\midrule
user       & user\_id, username,\newline email, password\_hash    & Authentication, authorship, and soft-delete via \texttt{is\_active} \\
\rowcolor{lightgray}
repository & repo\_id, user\_id (FK),\newline current\_hash (FK)  & Connection config for target DB; \texttt{current\_hash} serves as the HEAD pointer \\
commits    & hash, repo\_id (FK),\newline parent\_hash, author\_id (FK) & One row per commit; forms the linear Merkle chain \\
\rowcolor{lightgray}
snapshots  & id, commit\_hash (FK),\newline table\_name, ddl\_json, rows\_json & Full schema + data snapshot per table per commit \\
\end{tabularx}
\vspace{0.5em}

\subsection{Automatic Diff Engine}

The diff engine reads two snapshots from \texttt{gitdb\_meta} and generates valid
MySQL SQL to transition between them:

\begin{itemize}[leftmargin=1.5em, itemsep=0.3em]
    \item \textbf{Schema Diffing:} AST-level comparison using \texttt{sqlglot}
    to identify added tables, dropped tables, and altered columns
    \item \textbf{Data Diffing:} Primary-key-anchored algorithm producing
    \texttt{INSERT}, \texttt{UPDATE}, \texttt{DELETE} statements in $O(n)$ time
\end{itemize}

\subsection{Merkle Hash Chain}

Each commit is identified by a SHA-256 hash derived from the serialised snapshot
content and the parent commit's hash:

\begin{tcolorbox}[commandbox, title=Hash Construction]
\begin{lstlisting}
hash = SHA256(serialise(schema, data) + parent_hash)
\end{lstlisting}
\end{tcolorbox}

This forms an immutable, tamper-evident chain. Since GitDB supports only a
\textbf{linear commit history} (no branching or merging), each commit has at most
one parent and at most one child --- the history is strictly a linked list, not a
DAG. Any modification to a historical snapshot immediately invalidates all
subsequent hashes, making unauthorised history rewriting detectable.

\subsection{Two-Phase Checkout for MySQL DDL Safety}

MySQL DDL statements issue implicit commits and cannot be rolled back. GitDB
handles this through a two-phase checkout strategy:

\begin{tcolorbox}[warningbox]
\textbf{MySQL DDL Note:} Unlike SQLite (which supports transactional DDL),
MySQL's \texttt{ALTER TABLE}, \texttt{DROP TABLE}, and \texttt{CREATE TABLE}
statements auto-commit. A simple \texttt{BEGIN} / \texttt{ROLLBACK} cannot undo
them. GitDB's two-phase checkout accounts for this explicitly.
\end{tcolorbox}

\begin{tcolorbox}[infobox]
\textbf{Two-Phase Checkout Workflow:}
\begin{enumerate}[leftmargin=1.5em, itemsep=0.2em]
    \item \textbf{Phase 1 --- Schema:} Validate and apply all DDL statements
    (\texttt{ALTER TABLE}, \texttt{CREATE TABLE}, \texttt{DROP TABLE}). On
    failure, re-apply the original schema from the \texttt{gitdb\_meta} snapshot.
    \item \textbf{Phase 2 --- Data:} Apply all \texttt{INSERT} /
    \texttt{UPDATE} / \texttt{DELETE} statements inside an explicit
    \texttt{BEGIN} / \texttt{COMMIT} transaction. On failure, \texttt{ROLLBACK}
    is called automatically.
\end{enumerate}
Since the rollback target is always a known-good snapshot stored in
\texttt{gitdb\_meta}, the database is guaranteed to return to a consistent state
even when DDL auto-commit prevents traditional rollback. On success,
\texttt{repository.current\_hash} is updated to the checked-out commit hash.
\end{tcolorbox}

\subsection{Dual Interface Architecture}

\begin{itemize}[leftmargin=1.5em, itemsep=0.3em]
    \item \textbf{CLI:} Git-familiar commands (\texttt{init}, \texttt{commit},
    \texttt{log}, \texttt{diff}, \texttt{checkout}, \texttt{status}) built with
    \texttt{click} and \texttt{rich}
    \item \textbf{Web UI:} Visual commit history, side-by-side diff viewer,
    schema browser, built in React~18 + Tailwind CSS
\end{itemize}

% ── Section 3: System Architecture ───────────────────────────────────────────
\section{System Architecture}

\subsection{Overall Architecture}

GitDB consists of four loosely coupled layers:

\begin{enumerate}[leftmargin=1.5em, itemsep=0.3em]
    \item \textbf{Target MySQL DB} --- the database being versioned (any user
    MySQL database)
    \item \textbf{Core Engine} --- Python module handling snapshotting, hashing,
    diffing, and checkout
    \item \textbf{gitdb\_meta} --- the dedicated MySQL metadata database storing
    all user, repository, commit, and snapshot data
    \item \textbf{Interfaces} --- CLI (click/rich) and Web UI (React + Flask
    REST API)
\end{enumerate}

\subsection{Schema Introspection via \texttt{information\_schema}}

GitDB uses MySQL's \texttt{information\_schema} to programmatically capture the
full database schema at commit time:

\begin{tcolorbox}[commandbox, title=Schema Capture Queries]
\begin{lstlisting}[language=SQL]
-- Get all user tables
SELECT TABLE_NAME
FROM information_schema.TABLES
WHERE TABLE_SCHEMA = %s AND TABLE_TYPE = 'BASE TABLE';

-- Get column definitions per table
SELECT COLUMN_NAME, COLUMN_TYPE, IS_NULLABLE,
       COLUMN_DEFAULT, COLUMN_KEY, EXTRA
FROM information_schema.COLUMNS
WHERE TABLE_SCHEMA = %s AND TABLE_NAME = %s
ORDER BY ORDINAL_POSITION;

-- Get full DDL for AST parsing
SHOW CREATE TABLE `table_name`;
\end{lstlisting}
\end{tcolorbox}

\subsection{Diff Engine Design}

\subsubsection{Schema Diffing (AST Level)}

Schema comparison uses \texttt{sqlglot} to parse the captured
\texttt{SHOW CREATE TABLE} DDL into Abstract Syntax Trees. This allows GitDB to
understand SQL semantics rather than doing naive text comparison:

\begin{itemize}[leftmargin=1.5em, itemsep=0.3em]
    \item Added tables $\rightarrow$ generate \texttt{CREATE TABLE} statements
    \item Dropped tables $\rightarrow$ generate \texttt{DROP TABLE} statements
    \item Added columns $\rightarrow$ generate \texttt{ALTER TABLE ... ADD COLUMN}
    \item Dropped columns $\rightarrow$ generate \texttt{ALTER TABLE ... DROP COLUMN}
    \item Modified column definitions $\rightarrow$ generate
    \texttt{ALTER TABLE ... MODIFY COLUMN}
\end{itemize}

\subsubsection{Data Diffing (Primary-Key Anchored)}

Data comparison indexes both snapshots by primary key and compares them as
dictionaries:

\begin{enumerate}[leftmargin=1.5em, itemsep=0.3em]
    \item Both old and new snapshots are loaded from \texttt{gitdb\_meta.snapshots}
    and indexed by primary key
    \item Keys present in new but not old $\rightarrow$ \texttt{INSERT} statements
    \item Keys present in both but with differing values $\rightarrow$
    \texttt{UPDATE} statements
    \item Keys present in old but not new $\rightarrow$ \texttt{DELETE} statements
\end{enumerate}

This algorithm operates in $\mathbf{O(n)}$ time with respect to the number of
rows. Tables without a primary key cannot be diffed and are flagged at commit
time with a clear diagnostic message.

\subsection{\texttt{gitdb\_meta} Schema}

The complete DDL for the four-entity \texttt{gitdb\_meta} database is as follows.
Note that the \texttt{ALTER TABLE} on \texttt{repository} is required after
\texttt{commits} is created, to resolve the circular foreign key dependency
between the two tables.

\begin{tcolorbox}[commandbox, title=gitdb\_meta Complete DDL]
\begin{lstlisting}[language=SQL]
CREATE DATABASE IF NOT EXISTS gitdb_meta;
USE gitdb_meta;

-- Users (authentication and authorship)
CREATE TABLE user (
    user_id       INT          NOT NULL AUTO_INCREMENT,
    username      VARCHAR(50)  NOT NULL,
    email         VARCHAR(100) NOT NULL,
    password_hash VARCHAR(255) NOT NULL,
    full_name     VARCHAR(100) NOT NULL,
    is_active     BOOLEAN      NOT NULL DEFAULT TRUE,
    created_at    DATETIME     NOT NULL DEFAULT CURRENT_TIMESTAMP,
    PRIMARY KEY (user_id),
    UNIQUE KEY uq_username (username),
    UNIQUE KEY uq_email    (email)
);

-- Repositories (connection config + HEAD pointer)
CREATE TABLE repository (
    repo_id          INT          NOT NULL AUTO_INCREMENT,
    user_id          INT          NOT NULL,
    repo_name        VARCHAR(100) NOT NULL,
    target_db_name   VARCHAR(100) NOT NULL,
    db_host          VARCHAR(100) NOT NULL,
    db_port          INT          NOT NULL DEFAULT 3306,
    db_user          VARCHAR(50)  NOT NULL,
    db_password_key  VARCHAR(100) NULL,     -- keyring key, not the password
    current_hash     CHAR(64)     NULL,
    created_at       DATETIME     NOT NULL DEFAULT CURRENT_TIMESTAMP,
    PRIMARY KEY (repo_id),
    CONSTRAINT fk_repo_user
        FOREIGN KEY (user_id) REFERENCES user (user_id)
);

-- Commits (linear Merkle chain, no branching)
CREATE TABLE commits (
    hash         CHAR(64)     NOT NULL,
    repo_id      INT          NOT NULL,
    parent_hash  CHAR(64)     NULL,
    author_id    INT          NOT NULL,
    message      TEXT         NOT NULL,
    created_at   DATETIME     NOT NULL DEFAULT CURRENT_TIMESTAMP,
    PRIMARY KEY (hash),
    UNIQUE KEY uq_parent_hash (parent_hash),   -- enforces linear chain at DB level
    CONSTRAINT fk_commit_repo
        FOREIGN KEY (repo_id)     REFERENCES repository (repo_id),
    CONSTRAINT fk_commit_parent
        FOREIGN KEY (parent_hash) REFERENCES commits    (hash),
    CONSTRAINT fk_commit_author
        FOREIGN KEY (author_id)   REFERENCES user       (user_id)
);

-- Add HEAD FK on repository after commits table exists
ALTER TABLE repository
    ADD CONSTRAINT fk_repo_head
        FOREIGN KEY (current_hash) REFERENCES commits (hash);

-- Snapshots (full schema + data per table per commit)
CREATE TABLE snapshots (
    id           BIGINT       NOT NULL AUTO_INCREMENT,
    commit_hash  CHAR(64)     NOT NULL,
    table_name   VARCHAR(255) NOT NULL,
    ddl_json     LONGTEXT     NOT NULL,
    rows_json    LONGTEXT     NOT NULL,
    row_count    INT          NOT NULL,
    captured_at  DATETIME     NOT NULL DEFAULT CURRENT_TIMESTAMP,
    PRIMARY KEY (id),
    UNIQUE KEY uq_commit_table (commit_hash, table_name),
    CONSTRAINT fk_snap_commit
        FOREIGN KEY (commit_hash) REFERENCES commits (hash)
);
\end{lstlisting}
\end{tcolorbox}

% ── Section 4: Database Structure ────────────────────────────────────────────
\section{Database Structure and Entity Relationships}

\subsection{Entity Overview}

The \texttt{gitdb\_meta} database consists of four entities. The HEAD concept
from the original Git analogy is represented as \texttt{current\_hash} directly
inside \texttt{repository} rather than as a separate table, keeping the schema
lean and referentially clean.

\begin{center}
\renewcommand{\arraystretch}{1.4}
\begin{tabular}{>{\bfseries\ttfamily}l c c l}
\toprule
\textbf{Entity} & \textbf{PK Type} & \textbf{Rows (typical)} & \textbf{Storage concern} \\
\midrule
user       & INT AUTO\_INCREMENT  & Few               & Negligible \\
repository & INT AUTO\_INCREMENT  & One per DB        & Negligible \\
commits    & CHAR(64) SHA-256     & One per commit    & Small \\
snapshots  & BIGINT AUTO\_INCREMENT & tables $\times$ commits & Large (LONGTEXT blobs) \\
\bottomrule
\end{tabular}
\end{center}

\subsection{Entity Attribute Tables}

\subsubsection{USER}

\begin{center}
\renewcommand{\arraystretch}{1.3}
\begin{tabular}{>{\ttfamily}l >{\ttfamily}l l p{5cm}}
\toprule
\textbf{Attribute} & \textbf{Type} & \textbf{Constraint} & \textbf{Description} \\
\midrule
user\_id       & INT          & PK, AUTO\_INCREMENT  & Surrogate primary key \\
username       & VARCHAR(50)  & UNIQUE, NOT NULL     & Unique login handle \\
email          & VARCHAR(100) & UNIQUE, NOT NULL     & Unique email address \\
password\_hash & VARCHAR(255) & NOT NULL             & bcrypt / argon2 hash \\
full\_name     & VARCHAR(100) & NOT NULL             & Display name used in \texttt{gitdb log} \\
is\_active     & BOOLEAN      & DEFAULT TRUE         & Soft-delete flag \\
created\_at    & DATETIME     & DEFAULT CURRENT\_TIMESTAMP & Account creation timestamp \\
\bottomrule
\end{tabular}
\end{center}

\subsubsection{REPOSITORY}

\begin{center}
\renewcommand{\arraystretch}{1.3}
\begin{tabular}{>{\ttfamily}l >{\ttfamily}l l p{4.5cm}}
\toprule
\textbf{Attribute} & \textbf{Type} & \textbf{Constraint} & \textbf{Description} \\
\midrule
repo\_id         & INT          & PK, AUTO\_INCREMENT       & Surrogate primary key \\
user\_id         & INT          & FK $\to$ user, NOT NULL   & Owner of this repository \\
repo\_name       & VARCHAR(100) & NOT NULL                  & Human-readable name \\
target\_db\_name & VARCHAR(100) & NOT NULL                  & MySQL database being versioned \\
db\_host         & VARCHAR(100) & NOT NULL                  & Server hostname or IP \\
db\_port         & INT          & NOT NULL, DEFAULT 3306    & Server port \\
db\_user         & VARCHAR(50)  & NOT NULL                  & MySQL login username \\
db\_password\_key & VARCHAR(100) & NULL                      & OS keyring service key; password never stored here \\
current\_hash    & CHAR(64)     & FK $\to$ commits, NULL    & HEAD pointer; NULL before first commit \\
created\_at      & DATETIME     & DEFAULT CURRENT\_TIMESTAMP & Registration timestamp \\
\bottomrule
\end{tabular}
\end{center}

\begin{tcolorbox}[infobox]
\textbf{Note on \texttt{current\_hash}:} This column replaces a separate
\texttt{head} table. It is updated on every successful \texttt{gitdb commit}
and \texttt{gitdb checkout}. The MySQL connection password is \emph{never} stored
in \texttt{gitdb\_meta} or \texttt{.gitdb/config.json} --- it is stored in the
OS-level keyring (via the \texttt{keyring} library); only the keyring service key
is persisted in \texttt{db\_password\_key}.
\end{tcolorbox}

\subsubsection{COMMITS}

\begin{center}
\renewcommand{\arraystretch}{1.3}
\begin{tabular}{>{\ttfamily}l >{\ttfamily}l l p{4.5cm}}
\toprule
\textbf{Attribute} & \textbf{Type} & \textbf{Constraint} & \textbf{Description} \\
\midrule
hash         & CHAR(64)  & PK                              & SHA-256 content hash \\
repo\_id     & INT       & FK $\to$ repository, NOT NULL   & Repository this commit belongs to \\
parent\_hash & CHAR(64)  & FK $\to$ commits, NULL          & Parent commit; NULL for root commit \\
author\_id   & INT       & FK $\to$ user, NOT NULL         & User who authored this commit \\
message      & TEXT      & NOT NULL                        & Commit description \\
created\_at  & DATETIME  & DEFAULT CURRENT\_TIMESTAMP     & Commit timestamp \\
\bottomrule
\end{tabular}
\end{center}

\begin{tcolorbox}[warningbox]
\textbf{Linear chain constraint:} GitDB does not support branching. Therefore,
\texttt{parent\_hash} is unique across the table --- no two commits share the
same parent. The self-referencing foreign key is \texttt{NULL} only for the root
commit (the very first commit in a repository). The hash is computed as:
\[
\texttt{hash} = \text{SHA-256}\bigl(\,\text{serialise}(\textit{schema},\,
\textit{data}) \;\|\; \texttt{parent\_hash}\,\bigr)
\]
\end{tcolorbox}

\subsubsection{SNAPSHOTS}

\begin{center}
\renewcommand{\arraystretch}{1.3}
\begin{tabular}{>{\ttfamily}l >{\ttfamily}l l p{4.5cm}}
\toprule
\textbf{Attribute} & \textbf{Type} & \textbf{Constraint} & \textbf{Description} \\
\midrule
id           & BIGINT       & PK, AUTO\_INCREMENT             & Surrogate primary key \\
commit\_hash & CHAR(64)     & FK $\to$ commits, NOT NULL      & Commit this snapshot belongs to \\
table\_name  & VARCHAR(255) & NOT NULL                        & Name of the MySQL table snapshotted \\
ddl\_json    & LONGTEXT     & NOT NULL                        & Serialised DDL (schema) as JSON \\
rows\_json   & LONGTEXT     & NOT NULL                        & Serialised table rows as JSON array \\
row\_count   & INT          & NOT NULL                        & Row count at snapshot time \\
captured\_at & DATETIME     & DEFAULT CURRENT\_TIMESTAMP     & Snapshot capture timestamp \\
\bottomrule
\end{tabular}
\end{center}

A \texttt{UNIQUE} constraint on \texttt{(commit\_hash, table\_name)} ensures no
duplicate snapshot exists for the same table within the same commit. A hard limit
of \textbf{10,000 rows per table} is enforced at commit time.

\subsection{Relationship Table}

\begin{center}
\renewcommand{\arraystretch}{1.5}
\begin{tabular}{l l l l p{4.5cm}}
\toprule
\textbf{From} & \textbf{To} & \textbf{Name} & \textbf{Cardinality} & \textbf{Description} \\
\midrule
user       & repository & \textit{owns}      & 1:N      & One user owns many repositories \\
user       & commits    & \textit{authors}   & 1:N      & One user authors many commits \\
repository & commits    & \textit{contains}  & 1:N      & One repository contains many commits \\
commits    & commits    & \textit{child of}  & 1:0..1   & Linear chain; each commit has at most one parent and at most one child \\
commits    & snapshots  & \textit{has}       & 1:N      & One commit has one snapshot per tracked table \\
repository & commits    & \textit{points to} & 1:1      & \texttt{current\_hash} is the HEAD pointer \\
\bottomrule
\end{tabular}
\end{center}

% ── Section 5: Feasibility ────────────────────────────────────────────────────
\section{Feasibility and Design Constraints}

\subsection{MySQL-Only Target}

GitDB is scoped exclusively to MySQL. This focus enables deeper MySQL-specific
optimisations and avoids the complexity of multi-database abstraction layers. The
driver used is \texttt{mysql-connector-python}, with connection parameters (host,
port, user, password, database) supplied at \texttt{gitdb init} time and stored
in a local \texttt{.gitdb/config.json}.

\subsection{Handling MySQL DDL Auto-Commit}

This is the primary engineering challenge that differentiates GitDB-for-MySQL
from a generic database versioning tool. The solution is the two-phase checkout
described in Section~2.5. The key insight is: since the rollback target is always
a previously committed, known-good snapshot stored in \texttt{gitdb\_meta}, DDL
auto-commit only means we must re-apply the old schema from the snapshot rather
than relying on \texttt{ROLLBACK} --- the end result is identical.

\subsection{Handling Known Hard Problems}

\subsubsection{Foreign Key Constraint Ordering}
MySQL enforces foreign key constraints during DML operations. GitDB disables
constraint checks for the duration of the data patch phase and re-enables them
after:

\begin{tcolorbox}[commandbox, title=Foreign Key Handling (MySQL)]
\begin{lstlisting}[language=SQL]
SET FOREIGN_KEY_CHECKS = 0;
-- ... apply all INSERT / UPDATE / DELETE statements ...
SET FOREIGN_KEY_CHECKS = 1;
\end{lstlisting}
\end{tcolorbox}

Since the target state is always a known-good snapshot, re-enabling checks after
the patch guarantees referential integrity of the final state.

\subsubsection{Column Type Changes and Table Renames}
Rather than attempting to infer developer intent in ambiguous cases, GitDB
explicitly flags these as unsupported operations with clear diagnostic messages.
This mirrors the design philosophy of production-grade tools such as Prisma and
Atlas, which prioritise correctness over automatic handling of ambiguous
transformations.

\subsection{Data Scale Constraints}

A row limit of \textbf{10,000 rows per table} is enforced at commit time. This
keeps snapshot sizes in \texttt{gitdb\_meta} manageable and ensures in-memory
diffing completes within practical time bounds for a mini-project scope.

\subsection{Correctness by Construction}

Within documented constraints, GitDB's diff engine is correct by construction:

\begin{itemize}[leftmargin=1.5em, itemsep=0.3em]
    \item Produces deterministic, reproducible SQL from algorithmic
    transformation of two known snapshots
    \item No reliance on external APIs or non-deterministic heuristics
    \item All transformations are pure functions of input snapshots stored in
    \texttt{gitdb\_meta}
    \item Two-phase checkout guarantees the database returns to a consistent
    state on any failure
\end{itemize}

% ── Section 6: Technology Stack ───────────────────────────────────────────────
\section{Technology Stack}

\renewcommand{\arraystretch}{1.4}
\begin{tabularx}{\textwidth}{>{\bfseries}p{3.2cm} >{\ttfamily}p{3.8cm} X}
\rowcolor{tableheader}
\textcolor{white}{\normalfont\bfseries Domain} &
\textcolor{white}{\normalfont\bfseries Technology} &
\textcolor{white}{\normalfont\bfseries Purpose} \\
\midrule
Core Engine       & Python 3.10+           & Snapshot generation, diff engine, and MySQL interaction \\
\rowcolor{lightgray}
Database Driver   & mysql-connector-python & Native MySQL connectivity for both target DB and gitdb\_meta \\
Schema Parsing    & sqlglot                & AST-level schema diff engine for DDL parsing \\
\rowcolor{lightgray}
Hashing           & hashlib (SHA-256)      & Merkle hash chain for commit identification \\
CLI Framework     & click, rich            & Git-style subcommands and coloured terminal output \\
\rowcolor{lightgray}
Backend API       & Flask                  & Lightweight REST API exposing engine operations over HTTP \\
Frontend UI       & React 18               & Component-based visual interface \\
\rowcolor{lightgray}
Styling           & Tailwind CSS           & Utility-first CSS framework \\
Commit History    & react-flow             & Renders the linear commit history visually \\
\rowcolor{lightgray}
Diff Viewer       & @git-diff-view/react   & Side-by-side schema and data comparison (React~18 compatible) \\
Serialisation     & JSON (Python stdlib)   & Schema and row data storage in gitdb\_meta.snapshots \\
\end{tabularx}

\begin{tcolorbox}[infobox]
\textbf{Note on \texttt{react-flow}:} The original design described the commit
graph as a Directed Acyclic Graph (DAG). Since GitDB supports only a linear
commit history with no branching or merging, \texttt{react-flow} is used to
render a simple linear chain of commit nodes rather than a full DAG.
\end{tcolorbox}

% ── Section 7: CLI Reference ──────────────────────────────────────────────────
\section{CLI Commands Reference}

The command-line interface closely follows Git conventions for immediate
familiarity.

\begin{tcolorbox}[commandbox, title=\texttt{gitdb register}]
\begin{lstlisting}
gitdb register
\end{lstlisting}
\textcolor{codetext}{\small Creates a user in \texttt{gitdb\_meta.user}
with an argon2-hashed password via interactive prompts. No plaintext password
is passed as a CLI argument.}
\end{tcolorbox}

\begin{tcolorbox}[commandbox, title=\texttt{gitdb init}]
\begin{lstlisting}
gitdb init --host localhost --port 3306 \
           --user root --password secret \
           --database myapp_db \
           --repo-name myapp --author username
\end{lstlisting}
\textcolor{codetext}{\small Connects to the specified MySQL server, creates the
\texttt{gitdb\_meta} database and its four tables (\texttt{user},
\texttt{repository}, \texttt{commits}, \texttt{snapshots}) if they do not already
exist, registers the target database as a new repository, writes connection
config to \texttt{.gitdb/config.json}, and stores DB credentials using the
OS keyring service.}
\end{tcolorbox}

\begin{tcolorbox}[commandbox, title=\texttt{gitdb commit}]
\begin{lstlisting}
gitdb commit -m "<message>" --author USERNAME
\end{lstlisting}
\textcolor{codetext}{\small Introspects the target MySQL database via
\texttt{information\_schema} and \texttt{SHOW CREATE TABLE}, serialises the full
schema and data as JSON, computes the SHA-256 Merkle hash, and writes the commit
and snapshot records to \texttt{gitdb\_meta}. Updates \texttt{repository.current\_hash}
to the new commit. Resolves \texttt{--author} to an active user in
\texttt{gitdb\_meta.user}. Enforces the 10,000-row limit per table.}
\end{tcolorbox}

\begin{tcolorbox}[commandbox, title=\texttt{gitdb log}]
\begin{lstlisting}
gitdb log [--oneline] [--graph]
\end{lstlisting}
\textcolor{codetext}{\small Queries \texttt{gitdb\_meta.commits} and displays the
commit history in reverse chronological order by walking the \texttt{parent\_hash}
chain from \texttt{repository.current\_hash}. \texttt{--oneline} shows a compact
format; \texttt{--graph} renders an ASCII linear chain.}
\end{tcolorbox}

\begin{tcolorbox}[commandbox, title=\texttt{gitdb diff}]
\begin{lstlisting}
gitdb diff <hash1> <hash2> [--schema-only] [--data-only]
\end{lstlisting}
\textcolor{codetext}{\small Reads both snapshots from \texttt{gitdb\_meta.snapshots},
runs the diff engine, and outputs the MySQL SQL statements (\texttt{ALTER TABLE},
\texttt{CREATE TABLE}, \texttt{DROP TABLE}, \texttt{INSERT}, \texttt{UPDATE},
\texttt{DELETE}) required to transition from \texttt{hash1} to \texttt{hash2}.}
\end{tcolorbox}

\begin{tcolorbox}[commandbox, title=\texttt{gitdb checkout}]
\begin{lstlisting}
gitdb checkout <commit_hash>
\end{lstlisting}
\textcolor{codetext}{\small Executes a two-phase checkout against the target MySQL
database. Phase~1 applies all DDL changes; on failure, the prior schema is
restored from \texttt{gitdb\_meta}. Phase~2 applies all data changes inside an
explicit transaction with \texttt{SET FOREIGN\_KEY\_CHECKS = 0}; on failure, the
transaction is rolled back. On success, updates
\texttt{repository.current\_hash} to the checked-out commit hash.}
\end{tcolorbox}

\begin{tcolorbox}[commandbox, title=\texttt{gitdb status}]
\begin{lstlisting}
gitdb status
\end{lstlisting}
\textcolor{codetext}{\small Compares the live MySQL database state against the
HEAD snapshot referenced by \texttt{repository.current\_hash} and reports
uncommitted changes: added/dropped tables, modified columns, and row count
differences per table.}
\end{tcolorbox}

% ── Section 8: Web UI Architecture ───────────────────────────────────────────
\section{Web UI Architecture}

The React-based web interface provides visual tooling for exploring and managing
MySQL database versions.

\subsection{Commit History Visualisation}

Built with \texttt{react-flow}, the commit history is rendered as a linear
chain of nodes. Each node shows the commit hash, message, author
(\texttt{user.full\_name}), and timestamp. The current HEAD --- identified via
\texttt{repository.current\_hash} --- is visually highlighted.

\subsection{Schema and Data Diff Viewer}

Powered by \texttt{@git-diff-view/react}, the interface displays side-by-side
comparisons of DDL and row-level data changes between any two commits, with
syntax highlighting and a SQL patch preview before checkout.

\subsection{Schema Table Browser}

Provides a read-only view of the MySQL database schema at any commit, showing
table structure, column types, constraints, primary keys, and foreign keys as
recorded in \texttt{gitdb\_meta.snapshots}.

\subsection{Flask REST API}

The backend exposes GitDB engine operations over HTTP:

\begin{itemize}[leftmargin=1.5em, itemsep=0.3em]
    \item \texttt{GET /commits} --- List all commits for the current repository
    from \texttt{gitdb\_meta}
    \item \texttt{POST /commit} --- Create new commit, write snapshot, update
    \texttt{repository.current\_hash}
    \item \texttt{GET /diff/:hash1/:hash2} --- Compute and return the SQL diff
    \item \texttt{POST /checkout/:hash} --- Execute two-phase checkout against
    MySQL; update \texttt{repository.current\_hash} on success
    \item \texttt{GET /status} --- Report uncommitted changes vs HEAD snapshot
\end{itemize}

% ── Section 9: Academic Contributions ────────────────────────────────────────
\section{Academic Contributions and Learning Outcomes}

GitDB serves as a comprehensive applied demonstration of core DBMS concepts, all
exercised against MySQL specifically.

\subsection{MySQL Schema Introspection}

Implementation of programmatic schema access via \texttt{information\_schema.TABLES},
\texttt{information\_schema.COLUMNS}, \texttt{information\_schema.KEY\_COLUMN\_USAGE},
and \texttt{SHOW CREATE TABLE} demonstrates how a production RDBMS exposes its
own metadata and the structure of system catalogues.

\subsection{ACID Properties and MySQL DDL Constraints}

The two-phase checkout directly illustrates the ACID properties and their limits
in MySQL:

\begin{itemize}[leftmargin=1.5em, itemsep=0.3em]
    \item \textbf{Atomicity:} Data patch phase is all-or-nothing via explicit
    transaction
    \item \textbf{Consistency:} Referential integrity guaranteed through
    \texttt{FOREIGN\_KEY\_CHECKS} discipline
    \item \textbf{Isolation:} Transaction-level separation during data restore
    \item \textbf{Durability:} Committed state is permanent;
    \texttt{repository.current\_hash} reflects actual DB state
    \item \textbf{DDL Auto-Commit:} Demonstrates a real and practically
    important MySQL limitation and how to engineer around it
\end{itemize}

\subsection{Structural Schema Comparison}

The AST-level diff engine exercises concepts that appear in ORM migration
generation (Alembic, Django migrations), database CI/CD validation pipelines,
and schema synchronisation tools --- all of which are core to modern database
DevOps practice.

\subsection{Cryptographic Integrity Verification}

The SHA-256 Merkle hash chain demonstrates the application of cryptographic
hashing to build tamper-evident audit logs, a concept that generalises to
blockchain, distributed ledgers, and data integrity verification in enterprise
systems.

\subsection{Primary Key as Stable Row Identifier}

The primary-key-anchored data diff reinforces the foundational role of primary
keys in relational databases and illustrates the practical consequences of key
selection: tables without primary keys cannot be diffed and are flagged at
commit time with a clear diagnostic.

\subsection{Relational Design of the Versioning Layer}

The \texttt{gitdb\_meta} schema itself --- with its self-referential
\texttt{commits} table (parent hash foreign key), the normalised
\texttt{snapshots} table, the \texttt{user} table for authorship, and the
\texttt{current\_hash} HEAD pointer embedded in \texttt{repository} --- is a
teaching artefact in relational data modelling. It demonstrates how a linked-list
history structure, multi-user ownership, and a HEAD pointer can all be represented
cleanly in a four-table relational schema without redundancy.

% ── Section 10: Conclusion ────────────────────────────────────────────────────
\section{Conclusion}

GitDB represents a focused, well-scoped solution to a real and recognised problem
in MySQL database development: the absence of intuitive, Git-like version control
for database state. By adopting a state-based approach with a dedicated
\texttt{gitdb\_meta} metadata database, automatic diff generation, and a
two-phase checkout strategy that correctly handles MySQL's DDL auto-commit
behaviour, it provides MySQL developers with a familiar version control workflow
for managing both schema and data evolution.

The project's scope reflects deliberate engineering decisions:

\begin{itemize}[leftmargin=1.5em, itemsep=0.3em]
    \item \textbf{MySQL-only target} --- enables MySQL-specific optimisations
    and avoids multi-DB abstraction complexity
    \item \textbf{Four-entity \texttt{gitdb\_meta} schema} --- \texttt{user},
    \texttt{repository}, \texttt{commits}, \texttt{snapshots}; HEAD pointer
    embedded in \texttt{repository.current\_hash}
    \item \textbf{Two-phase checkout} --- principled handling of MySQL DDL
    auto-commit without pretending it does not exist
    \item \textbf{Linear commit history} --- no branching or merging, keeping
    scope tractable; history is a linked list not a DAG
    \item \textbf{10,000-row limit per table} --- documented constraint
    appropriate to mini-project scale
    \item \textbf{Column type changes and renames flagged, not silently handled}
    --- correctness over convenience, consistent with Prisma and Atlas
    \item \textbf{User entity for authorship and authentication} --- commits
    carry \texttt{author\_id} FK; \texttt{full\_name} used in log output
\end{itemize}

From an academic perspective, GitDB demonstrates substantive breadth across core
DBMS concepts --- MySQL schema introspection, ACID transaction properties, DDL
auto-commit constraints, structural AST-level diffing, cryptographic Merkle
hashing, and relational data modelling of the versioning layer itself --- while
providing each team member with ownership of a well-defined, independently
testable subsystem.

The result is a working system that brings Git's proven version control paradigm
to the MySQL layer, implemented with established algorithms, mature libraries,
and principled handling of MySQL's real-world constraints.

\end{document}